\begin{explanationbox}
\textbf{一、先从熟悉的圆出发：圆的切线为什么简单？}
\begin{enumerate}
\item 对圆 $x^2+y^2=R^2$，圆心 $O$ 到切线的垂足就是切点 $P$，所以 $OP=R$，且半径垂直切线。
\item 椭圆的焦点 $F_1,F_2$ 不在中心；看似每个焦点到不同切线的距离都会变，\textbf{但两个距离一增一减，乘积始终为 $b^2$}。
\item 更奇妙：把两个焦点分别向切线作垂线，垂足并不留在椭圆上，而是落在圆 $x^2+y^2=a^2$ 上。
这相当于把一个椭圆切线问题，转换成初中熟悉的\textbf{圆的中位线与点幂问题}。
\end{enumerate}

\textbf{二、为什么椭圆会有反射性质？先想一根长度 $2a$ 的绳子。}

把绳子两端固定在 $F_1,F_2$，用笔尖 $P$ 绷紧绳子：
$PF_1+PF_2=2a$。
笔尖在椭圆上挪动时，这个和始终不变。
在 $P$ 的切线方向挪动极小一段，靠近一个焦点而减少的长度，
必须恰好补偿远离另一焦点而增加的长度。
因此两条焦点连线和切线所成的\textbf{锐角相等}；
等价地，法线平分两焦点射线所成的内角。
这就是椭圆镜面把一个焦点的光反射到另一焦点的直观原因。

\textbf{三、由反射到共圆：只需要初中三角形中位线。}

把焦点 $F_1$ 沿切线 $l$ 对折到点 $Q$。
由于切线反射性质，$F_2,P,Q$ 共线，且 $P$ 在 $F_2,Q$ 之间；
于是
\[
F_2Q=F_2P+PQ=F_2P+PF_1=2a.
\]
反射轴 $l$ 垂直平分 $F_1Q$，所以 $H_1$ 是 $F_1Q$ 中点；
而 $O$ 是 $F_1F_2$ 中点。
在三角形 $F_1F_2Q$ 中，$OH_1$ 是中位线：
\[
OH_1=\frac{F_2Q}{2}=a.
\]
交换两焦点，同样得到 $OH_2=a$。
\textbf{记忆口令：焦点对折、长度 $2a$、中位线变 $a$。}

\textbf{四、由共圆到距离积：圆的点幂定理突然出现。}

令 $K$ 是主圆上与 $H_2$ 关于圆心 $O$ 对称的点（$K=-H_2$）。
由于 $F_1=-F_2$，且 $F_1H_1$ 和 $F_2H_2$ 都垂直于切线，
$H_1,F_1,K$ 共线；且 $F_1H_1=d_1$、$F_1K=d_2$。
点 $F_1$ 位于半径 $a$ 的主圆内部，应用\textbf{圆内一点的点幂}：
\[
F_1H_1\cdot F_1K=a^2-OF_1^2=a^2-c^2=b^2.
\]
也就是 $d_1d_2=b^2$。
\textbf{记忆口令：对径点 $-H_2$、过焦点同一直线、点幂得 $b^2$。}

\textbf{五、面积为什么需要分类？}

四边形 $F_1H_1H_2F_2$ 是以 $F_1H_1,F_2H_2$ 为平行边的梯形（竖直切线时退化），所以
\[
S=\frac{d_1+d_2}{2}\,H_1H_2.
\]
由 $F_1F_2=2c$ 的水平、垂直分解，得到
\[
H_1H_2^2+(d_1-d_2)^2=4c^2.
\]
再用 $d_1d_2=b^2$，可化为
\[
\boxed{H_1H_2^2+(d_1+d_2)^2=4a^2}.
\]
因此把 $H_1H_2$ 与 $d_1+d_2$ 看成直角三角形的两条直角边、$2a$ 看成斜边，
面积就是两直角边乘积的一半，故 $S\le a^2$。
但是\textbf{直角边相等的构型不一定能由给定椭圆的切线实现}：
当 $c<b$ 时，椭圆不够扁，实际最大值只有 $2bc$；
当 $c\ge b$ 时，才能达到上界 $a^2$。
这是\textbf{不等式求出的上界未必可以取到}的典型例子。

\textbf{六、截距积如何一眼记住？}

切线方程是 $\frac{x_0x}{a^2}+\frac{y_0y}{b^2}=1$，
令 $y=0$ 便得 $x_N=a^2/x_0$。
法线垂直于切线，方向向量可取
$\bigl(x_0/a^2,y_0/b^2\bigr)$；
从 $P$ 沿法线走到 $x$ 轴得
$x_M=c^2x_0/a^2$。
\textbf{两式的 $x_0$ 与 $a^2$ 正好约掉，留下 $c^2$。}

\textbf{七、跟圆有什么相同与不同？}
\begin{itemize}
\item \textbf{圆的切线}：圆心到切线的垂足就是切点；\textbf{椭圆的切线}：焦点到切线的垂足都落在半径 $a$ 的辅助圆上。
\item \textbf{圆的半径垂直切线}：椭圆则由\textbf{焦点角的内角平分线}充当法线。
\item \textbf{圆的点幂 $a^2-c^2$}：在椭圆中转化为两个焦点到切线距离的乘积。
\item 椭圆趋于圆（$c\to0$）时，两个焦点合并到圆心，$d_1d_2\to a^2$、$H_1,H_2\to P$；
焦点三角形及两个角平分线截距的表述在极限中退化，不能直接照搬。
\end{itemize}

\textbf{使用信号}：看见``焦点到切线的垂距''、``切线垂足''、``反射''、``法线/角平分线''、
``截距乘积''、``面积最大''，先检查能否沿着\textbf{反射—共圆—点幂}这条链处理。
\end{explanationbox}
